\documentclass[11pt,a4paper]{article}
\usepackage{algorithm}
\usepackage{algpseudocode}
\begin{document}
\section{Pseudocode}
Algorithm~\ref{alg:euclid} is Euclid's method, set with line numbers in the
\texttt{ruled} float that the \texttt{algorithm} package declares.

\begin{algorithm}[H]
\caption{Euclid's algorithm}\label{alg:euclid}
\begin{algorithmic}[1]
\Require integers $a \ge 0$ and $b \ge 0$
\Ensure the greatest common divisor of $a$ and $b$
\Procedure{Euclid}{$a, b$}\Comment{The g.c.d.\ of a and b}
  \State $r \gets a \bmod b$
  \While{$r \ne 0$}\Comment{We have the answer if r is 0}
    \State $a \gets b$
    \State $b \gets r$
    \State $r \gets a \bmod b$
  \EndWhile
  \State \Return $b$\Comment{The gcd is b}
\EndProcedure
\end{algorithmic}
\end{algorithm}

The text continues after the algorithm, long enough to reach a second line
of this paragraph so that its spacing can be measured.

\begin{algorithm}[H]
\caption{Sum of the positive entries}\label{alg:sum}
\begin{algorithmic}
\Function{PositiveSum}{$x_1, \dots, x_n$}
  \State $s \gets 0$
  \For{$i \gets 1$ to $n$}
    \If{$x_i > 0$}
      \State $s \gets s + x_i$
    \ElsIf{$x_i = 0$}
      \State skip it
    \Else
      \State count it as negative
    \EndIf
  \EndFor
  \Repeat
    \State $s \gets$ \Call{Halve}{$s$}
  \Until{$s < 1$}
  \State \Return $s$
\EndFunction
\end{algorithmic}
\end{algorithm}

After both algorithms the text continues, and Algorithm~\ref{alg:sum} has
no line numbers.
\end{document}
